%--------------------------------------------------------------------------
% ABSTRACT - DEGREE PROJECT
%--------------------------------------------------------------------------
\chapter*{Abstract}
\addcontentsline{toc}{chapter}{Abstract}

This degree project develops a modular web application for document management, operational traceability, and administrative closure follow-up of work orders at CERMONT S.A.S. The diagnosis is based on fragmentation across paper forms, spreadsheets, and closure documents, a condition that hinders operational history consultation, coordination between areas, and timely consolidation of supporting records.

The proposal integrates user and role management, work orders, resource planning, photographic evidence, document generation, and administrative follow-up. Its purpose is to transform scattered information into structured, reusable, and auditable records. The architecture follows a modular web approach, separating responsibilities between presentation, business logic, and document-oriented persistence, which supports maintainability, traceability, and access control.

As a complementary extension, the work proposes a future evolution toward dynamic forms generated from the company's legacy documents, provided that formal approval is granted for its inclusion within the project scope. Functional validation is organized through use scenarios, integration tests, and role-based acceptance. Quantitative indicators of operational impact must be reported only when verifiable evidence is included in the appendices.

\vspace{0.5cm}
\noindent\textbf{Keywords:} document management, operational traceability, work orders, administrative closure, digital forms, modular web architecture, CERMONT S.A.S.
